\documentclass[aspectratio=169]{beamer}

\usetheme{Madrid}
\usecolortheme{default}
\setbeamertemplate{navigation symbols}{}

\usepackage{amsmath, amssymb}

\title{Orbitwise MDL Selection for Relative Symmetry Repair}
\author{}
\date{\today}

\begin{document}

\begin{frame}
  \titlepage
\end{frame}

\begin{frame}{Claim in One Sentence}
  \begin{itemize}
    \item For a fixed relative-periodic candidate $(p,\mathbf{s})$, the fitting problem splits into independent orbit classes.
    \item On each orbit class, the optimal binary background value is chosen by majority vote.
    \item Model selection then reduces to balancing improved residual fit against a complexity penalty for the number of orbit classes.
    \item Over a finite candidate set, an MDL/BIC-type selector stabilizes once the empirical orbit statistics and residual codelength rates converge.
  \end{itemize}
\end{frame}

\begin{frame}{Setup}
  Let
  \[
    U_T \in \{0,1\}^{T \times D_1 \times \cdots \times D_d}
  \]
  be a binary spacetime prefix, and let a relative-periodic candidate be indexed by
  \[
    m = (p,\mathbf{s}), \qquad \mathbf{s} \in \mathbb{Z}^d, \quad p \ge 1.
  \]

  The model class consists of all backgrounds $B$ satisfying
  \[
    B[t+p,\mathbf{x}+\mathbf{s}] = B[t,\mathbf{x}]
  \]
  with spatial coordinates interpreted modulo $(D_1,\dots,D_d)$.

  \vspace{0.5em}
  This induces
  \[
    k(m) = p \prod_{i=1}^d D_i
  \]
  orbit classes, and $B$ is constant on each orbit class.
\end{frame}

\begin{frame}{Terminology I}
  \begin{description}
    \item[Relative-periodic model] A background obeying $B[t+p,\mathbf{x}+\mathbf{s}] = B[t,\mathbf{x}]$.
    \item[Orbit class] The set of sites linked by repeated application of the shift-and-advance map $(t,\mathbf{x}) \mapsto (t+p,\mathbf{x}+\mathbf{s})$.
    \item[Template parameter] The single binary value assigned to one orbit class.
    \item[Orbitwise projection] The nearest admissible background obtained by optimizing independently on each orbit class.
    \item[Majority vote] The Hamming-loss minimizer on an orbit class.
  \end{description}
\end{frame}

\begin{frame}{Terminology II}
  \begin{description}
    \item[Empirical orbit frequency] The observed fraction of $1$s in a given orbit class.
    \item[Defect mask] $M = U_T \oplus B^*$, marking disagreement between the data and the fitted background.
    \item[Residual codelength] The number of bits used to encode the defect mask under a chosen residual code.
    \item[Parametric penalty] A complexity term of order $\frac{k(m)}{2}\log T$.
    \item[Selector stabilizes] Beyond some observation length, the minimizing candidate no longer changes.
    \item[Dimension-agnostic] The argument depends only on the orbit partition, not on whether the spacetime is 1D, 2D, or 3D.
  \end{description}
\end{frame}

\begin{frame}{Theorem 1: Orbitwise Reduction}
  \begin{block}{Statement}
    Fix a candidate $m=(p,\mathbf{s})$ and let $\{O_j\}_{j=1}^{k(m)}$ be its orbit classes. For each class $O_j$, let
    \[
      n_j^{(1)} = \#\{(t,\mathbf{x}) \in O_j : U_T[t,\mathbf{x}] = 1\},
      \qquad
      n_j^{(0)} = |O_j| - n_j^{(1)}.
    \]
    Then the Hamming-optimal background $B^*_{m,T}$ is obtained by choosing, independently for each orbit class,
    \[
      b_j =
      \begin{cases}
        1 & \text{if } n_j^{(1)} > n_j^{(0)}, \\
        0 & \text{if } n_j^{(0)} > n_j^{(1)}, \\
        \text{either value} & \text{if } n_j^{(0)} = n_j^{(1)}.
      \end{cases}
    \]
  \end{block}

  \[
    d(U_T,B^*_{m,T}) = \sum_{j=1}^{k(m)} \min\bigl(n_j^{(0)}, n_j^{(1)}\bigr).
  \]
\end{frame}

\begin{frame}{Why This Looks Like Independent Bernoulli Estimation}
  \begin{itemize}
    \item Once the orbit partition is fixed, each orbit class contributes its own binary summary statistic.
    \item The optimal class label depends only on the empirical frequency within that class.
    \item There is no coupling between classes in the Hamming objective.
    \item So the global fitting problem reduces to $k(m)$ separate binary decisions.
  \end{itemize}

  \vspace{0.5em}
  Heuristically: each orbit class acts like a repeated sample from a class-specific Bernoulli source, estimated by majority vote.
\end{frame}

\begin{frame}{Selection Score}
  For each candidate $m$, define the defect mask
  \[
    M^*_{m,T} = U_T \oplus B^*_{m,T}.
  \]

  Let $L_T(M^*_{m,T})$ be the residual codelength. Consider the score
  \[
    S_T(m) = L_T(M^*_{m,T}) + \frac{k(m)}{2}\log T.
  \]

  \begin{itemize}
    \item The first term rewards better fit.
    \item The second term penalizes model capacity.
    \item Since $k(m) = p \prod_i D_i$, larger periods create more orbit classes and therefore higher complexity.
  \end{itemize}
\end{frame}

\begin{frame}{Theorem 2: Finite-Candidate Stabilization}
  \begin{block}{Assumptions}
    Let $\mathcal{M}$ be a finite set of candidates. Assume:
    \begin{enumerate}
      \item For each candidate $m$, the empirical $1$-frequency in every orbit class converges as $T \to \infty$.
      \item For each candidate $m$, the normalized residual codelength converges:
      \[
        \frac{L_T(M^*_{m,T})}{T} \to \ell(m).
      \]
      \item The limiting residual rate $\ell(m)$ has a unique minimizer $m_0$ over $\mathcal{M}$.
    \end{enumerate}
  \end{block}

  Then the selector
  \[
    \hat m_T = \arg\min_{m \in \mathcal{M}} S_T(m)
  \]
  stabilizes to $m_0$ for all sufficiently large $T$.
\end{frame}

\begin{frame}{Interpretation of Theorem 2}
  \begin{itemize}
    \item If two models have different asymptotic residual rates, the $O(T)$ residual term dominates.
    \item If several models tie in asymptotic residual rate, the $O(\log T)$ complexity term breaks the tie in favor of smaller $k(m)$.
    \item This is the sense in which over-parameterized refinements stop winning: they cannot justify their extra degrees of freedom unless they improve the asymptotic residual rate.
  \end{itemize}

  \vspace{0.5em}
  So the selector eventually settles on the simplest candidate with the best limiting residual description.
\end{frame}

\begin{frame}{Proof Sketch}
  \begin{enumerate}
    \item Orbitwise reduction makes the optimal fit explicit for every candidate.
    \item Convergence of empirical orbit frequencies implies that each candidate has a well-defined limiting defect behavior.
    \item Convergence of normalized residual codelength gives a limiting residual rate $\ell(m)$.
    \item Because the candidate set is finite, score differences can be compared pairwise.
    \item Residual-rate gaps scale like $T$.
    \item Complexity gaps scale like $\log T$.
    \item Hence the minimizing candidate eventually stops changing.
  \end{enumerate}
\end{frame}

\begin{frame}{What Can Safely Be Claimed}
  \begin{itemize}
    \item The orbitwise projection theorem is exact.
    \item The stabilization result is valid over a finite candidate set under explicit convergence assumptions.
    \item The argument is dimension-agnostic.
  \end{itemize}

  \begin{alertblock}{Important wording}
    If the residual term is a heuristic code such as run-length encoding, the safest phrase is:
    \begin{center}
      \emph{``MDL/BIC-type finite-candidate stabilization''}
    \end{center}
    rather than full classical BIC consistency.
  \end{alertblock}
\end{frame}

\begin{frame}{What Should Not Be Overclaimed}
  \begin{itemize}
    \item If defects themselves have genuine periodic structure, a larger period can absorb them into the template.
    \item In that case, the larger model may improve the asymptotic residual rate by an $O(T)$ amount.
    \item Then MDL should prefer the larger period, because it really describes the full spacetime better.
    \item So the selector is converging to the best description of the observed spacetime, not necessarily to a metaphysically separate ``background period''.
  \end{itemize}
\end{frame}

\begin{frame}{Analogy}
  Imagine a drifting wallpaper pattern in a looping video with occasional paint splatters.

  \begin{itemize}
    \item An orbit class is all pixels that correspond to the same wallpaper location after accounting for the drift and loop period.
    \item To recover the wallpaper, stack those pixels together and take a majority vote on their color.
    \item A more complicated wallpaper template can absorb more splatters.
    \item MDL charges ``rent'' for every extra tile in that template.
  \end{itemize}

  With enough frames, fake complexity stops being worth the rent.
\end{frame}

\begin{frame}{Role of the Experiments}
  The 1D, 2D, and 3D experiments do not prove the theorem.

  They support the assumptions by showing that:
  \begin{itemize}
    \item empirical orbit statistics appear to settle over longer observation windows,
    \item MDL-optimal periods stop drifting after sufficient data,
    \item the same reasoning works across structurally different dynamics.
  \end{itemize}

  So the experiments are validation of plausibility and scope, not substitutes for the proof.
\end{frame}

\begin{frame}{Takeaway}
  \begin{itemize}
    \item Relative-periodic fitting is orbitwise and explicit.
    \item Model selection is a competition between residual description and orbit-count complexity.
    \item Over finite candidate sets, stabilization follows once orbit statistics and residual codelength rates converge.
    \item The theorem is dimension-agnostic because it only uses the orbit partition.
  \end{itemize}
\end{frame}

\end{document}
